\documentclass{report}

\usepackage[utf8]{inputenc}
\usepackage[spanish]{babel}
\usepackage{geometry}
\usepackage{graphicx}
\usepackage{titlesec}
\usepackage{lipsum}
\usepackage[dvipsnames]{xcolor}
\usepackage[fleqn]{mathtools}
\usepackage{booktabs}
\usepackage{amsmath}
\usepackage{amssymb} % Para los símbolos de conjuntos de números
\usepackage{latexsym}
\usepackage{nccmath}
\usepackage{multicol}
\usepackage{listings}
\usepackage{tasks}
\usepackage{color}
\usepackage{float}
\usepackage{enumitem}
\usepackage{longtable}
\usepackage{makecell}
\usepackage{caption}
\usepackage[parfill]{parskip}
\usepackage{lipsum}
\usepackage{enumitem}
\usepackage[hidelinks]{hyperref}% Para manejar enlaces
\usepackage{fancyhdr}
\usepackage{chappg}
\usepackage{adjustbox}
\usepackage{tocloft}
\usepackage{soul}
\usepackage{tikz}
\usetikzlibrary{arrows}
\usepackage{pgfplots}
\pgfplotsset{compat=1.18}
%Definicion de colores
\definecolor{colorIPN}{rgb}{0.5, 0.0,0.13}
\definecolor{colorESCOM}{rgb}{0.0, 0.5,1.0}
\definecolor{myorange}{RGB}{255, 204, 153}
\graphicspath{ {imagenes/} }
%Configuraciones extras de colores
\definecolor{codegreen}{rgb}{0,0.6,0}
\definecolor{codered}{rgb}{0.6,0,0}
\definecolor{codegray}{rgb}{0.5,0.5,0.5}
\definecolor{codepurple}{rgb}{0.58,0,0.82}
\definecolor{backcolour}{rgb}{0.95,0.95,0.92}

% Configuración de la portada
\titleformat{\section}[block]{\normalfont\huge\bfseries}{\thesection}{1em}{}
\geometry{a4paper, margin=1in}

% Definir comandos para información repetitiva
\newcommand{\logoInstitucion}{recursos/logotipo_ipn.png} % Reemplaza con el nombre del archivo de tu logo
\newcommand{\logoUniversidad}{recursos/EscudoESCOM.png} % Reemplaza con el nombre del archivo de tu logo
\newcommand{\nombreInstituto}{Instituto Politecnico Nacional}
\newcommand{\facultad}{Escuela Superior de Computo}
\newcommand{\materia}{Tópicos Selectos de Algoritmos Bioinspirados}
\newcommand{\grupo}{7BM1}
\newcommand{\profesora}{Daniel Molina Pérez}
\newcommand{\periodo}{2025/02}

\newcommand{\alumnoA}{Carrillo Barreiro José Emiliano}
\newcommand{\alumnoB}{Martinez Ayala Gerardo}
\newcommand{\alumnoC}{Robles Otero José Ángel}
\newcommand{\alumnoD}{Vásquez Morales Haniel Ulises}
\newcommand{\alumnos}{Carrillo J., Martinex G., Robles J. \& Vásquez H.}

% Definir nuevas listas para información repetitiva
\newlist{bracketed}{enumerate}{1}
\setlist[bracketed,1]{label={[{\arabic*}]},left=0pt}

% Configuración del encabezado y pie de página
\pagestyle{fancy}
\fancyhf{} % Limpia todos los campos

% Encabezado
\fancyhead[C]{\textbf{\nombreInstituto\\ \facultad}}
\fancyhead[L]{\includegraphics[width=1cm]{recursos/logotipo_ipn.png}}
\fancyhead[R]{\includegraphics[width=1.5cm]{recursos/EscudoESCOM.png}}

% Pie de página
\fancyfoot[C]{\materia||\grupo\\ \alumnos\\ \thepage}
% Configuración del grosor de la línea en el encabezado y pie de página
\renewcommand{\headrulewidth}{0.4pt}
\renewcommand{\footrulewidth}{0.4pt}

% Ajuste del tamaño del encabezado
\setlength{\headheight}{34.0845pt}
\addtolength{\topmargin}{-22.0845pt}

\fancypagestyle{plain}{%
    \fancyhf{} % Borra encabezado y pie de página por defecto
    \fancyhead[C]{\textbf{\nombreInstituto\\ \facultad}}
    \fancyhead[L]{\includegraphics[width=1cm]{recursos/logotipo_ipn.png}}
    \fancyhead[R]{\includegraphics[width=1.5cm]{recursos/EscudoESCOM.png}}
    \fancyfoot[C]{\materia||\grupo\\ \alumnos\\ \thepage}
    \renewcommand{\headrulewidth}{0.4pt}
    \renewcommand{\footrulewidth}{0.4pt}
}


% Configuración de lstlisting para Python
\lstdefinestyle{mystyle}{
    backgroundcolor=\color{backcolour},
    commentstyle=\color{codegreen},
    keywordstyle=\color{codepurple},
    numberstyle=\tiny\color{codegray},
    stringstyle=\color{codered},
    basicstyle=\ttfamily\small,
    breakatwhitespace=false,
    breaklines=true,
    captionpos=b,
    keepspaces=true,
    numbers=left,
    numbersep=5pt,
    showspaces=false,
    showstringspaces=false,
    showtabs=false,
    tabsize=2
}

\lstdefinestyle{custompython}{
    language=Python,
    basicstyle=\small\ttfamily,
    keywordstyle=\color{blue},
    stringstyle=\color{green},
    commentstyle=\color{red},
    numbers=left,
    numberstyle=\tiny\color{gray},
    stepnumber=1,
    numbersep=5pt,
    backgroundcolor=\color{white},
    frame=single,
    captionpos=b,
    breaklines=true,
    breakatwhitespace=false,
    showstringspaces=false,
    morekeywords={import,plt,np,ws,y,gradient}
}

% Definición de estilo para C++
\lstdefinestyle{cppstyle}{
    language=C++,
    basicstyle=\small\ttfamily,
    keywordstyle=\color{blue},
    commentstyle=\color{green!40!black},
    stringstyle=\color{orange},
    showstringspaces=false,
    breaklines=true,
    breakatwhitespace=true,
    numbers=left,
    numberstyle=\tiny,
    stepnumber=1,
    numbersep=5pt,
    frame=single,
    captionpos=b,
    backgroundcolor=\color{gray!5},
    xleftmargin=0.5cm,
    xrightmargin=0.5cm
}

% Definición de estilo para MATLAB
\lstdefinestyle{matlabstyle}{
    language=Matlab,
    basicstyle=\small\ttfamily,
    keywordstyle=\color{blue},
    commentstyle=\color{green!40!black},
    stringstyle=\color{purple!50!blue!50!white},
    showstringspaces=false,
    breaklines=true,
    breakatwhitespace=true,
    numbers=left,
    numberstyle=\tiny,
    stepnumber=1,
    numbersep=5pt,
    frame=single,
    captionpos=b,
    backgroundcolor=\color{gray!5},
    xleftmargin=0.5cm,
    xrightmargin=0.5cm
}

%Definicion de estilo para Python
\lstdefinestyle{pythonstyle}{
    language=Python,
    basicstyle=\small\ttfamily,
    keywordstyle=\color{red!80!black},
    commentstyle=\color{gray!40!white},
    stringstyle=\color{yellow!30!orange!80!white},
    showstringspaces=false,
    breaklines=true,
    breakatwhitespace=true,
    numbers=left,
    numberstyle=\tiny,
    stepnumber=1,
    numbersep=5pt,
    frame=single,
    captionpos=b,
    backgroundcolor=\color{gray!5},
    xleftmargin=0.5cm,
    xrightmargin=0.5cm
}

% Personalización del índice de listados
\newcommand{\listcppname}{Índice de Listados de C++}
\newcommand{\listmatlabname}{Índice de Listados de MATLAB}
\newcommand{\listpythonname}{Índice de Listados de Python}
\newlistof{cpplistings}{lolcpp}{\listcppname}
\newlistof{matlablistings}{lolmatlab}{\listmatlabname}
\newlistof{pythonlistings}{lolpython}{\listpythonname}

% Configuración para los listados de C++
\lstset{style=cppstyle}
\newcommand{\cpplisting}[1]{%
    \refstepcounter{cpplistings}%
    \addcontentsline{lolcpp}{cpplistings}{\protect\numberline{\thecpplistings}#1}\par}

% Configuración para los listados de MATLAB
\lstset{style=matlabstyle}
\newcommand{\matlablisting}[1]{%
    \refstepcounter{matlablistings}%
    \addcontentsline{lolmatlab}{matlablistings}{\protect\numberline{\thematlablistings}#1}\par}

% Configuración para los listados de Python
\lstset{style=pythonstyle}
\newcommand{\pythonlisting}[1]{%
    \refstepcounter{pythonlistings}%
    \addcontentsline{lolpython}{pythonlistings}{\protect\numberline{\thepythonlistings}#1}\par}

\begin{document}

% Portada
% portada
\begin{titlepage}
    \begin{center}
        \vspace*{1cm}

        \begin{tabular}{c@{\hspace{2cm}}c}
            \includegraphics[width=0.25\textwidth]{\logoInstitucion} &
            \includegraphics[width=0.4\textwidth]{\logoUniversidad}
        \end{tabular}

        \vspace{1cm}

        \textbf{\LARGE \nombreInstituto} \\
        \textbf{\Large \facultad} \\
        \vspace{0.5cm}
        \textbf{\large Materia: \materia} \\
        \textbf{\large Grupo: \grupo} \\
        \vspace{0.5cm}
        \textbf{\large Profesor: \profesora} \\
        \textbf{\large Periodo: \periodo} \\

        \vspace{0.75cm}

        \textbf{\LARGE Practica 01} \\
        \vspace{0.5cm}
        \textbf{\Large \textit{Maximizar Contraste en Imagenes Medicas.}} \\

        \vspace{0.3cm}

        \textbf{\large Realizado por:} \\
        \textbf{\large \alumnoA \\ \alumnoB \\ \alumnoC \\ \alumnoD}

        % Bloque de Resumen y Abstract
        \begin{minipage}{0.8\textwidth}
            \textbf{Abstract:}\\[0.3cm]
            Medical image processing is key for diagnosis, where contrast enhancement is essential for the correct identification of structures. In this practice, a bio-inspired algorithm is implemented to optimize metrics such as standard deviation (STD) or entropy to improve visual quality without distorting the original information. Using techniques based on swarm intelligence and natural evolution, the method dynamically adjusts the contrast enhancement parameters. The results demonstrate better visual perception and greater discrimination of key anatomical structures, facilitating medical interpretation.

        \end{minipage}

        \begin{minipage}{0.8\textwidth}
            \textbf{Resumen:}\\[0.3cm]
            Las imágenes médicas suelen presentar bajo contraste, lo que dificulta su análisis y diagnóstico. Para abordar este problema, se propone un algoritmo bioinspirado que optimiza métricas como la desviación estándar y la entropía, permitiendo mejorar la calidad visual sin comprometer la información original. Este enfoque adaptativo maximiza el detalle de las estructuras relevantes, facilitando su interpretación y apoyando la toma de decisiones clínicas.
        \end{minipage}
        
        \vspace{0.3cm}

        \textbf{\large Fecha: \today}

    \end{center}
\end{titlepage}


% Índice
\tableofcontents
\newpage
%\listofcpplistings%Crea indice de codigos en C++
%\listofmatlablistings%Crea indice de codigos en MATLAB
\listoftables
\listoffigures
%\listoftables

%Capitulos
\chapter{Introducción}
El análisis de imágenes médicas desempeña un papel fundamental en el diagnóstico, la planificación del tratamiento y el seguimiento de diversas condiciones médicas. Dentro de este campo, el mejoramiento del contraste se erige como una etapa crucial para optimizar la calidad visual de las imágenes, facilitando la identificación de anomalías y la extracción de características relevantes. Las imágenes médicas a menudo presentan desafíos inherentes, como bajo contraste, ruido, pérdida de señal y estructuras anatómicas complejas, lo que puede dificultar un análisis y diagnóstico precisos. En este contexto, las exploraciones por tomografía computarizada (TC) con contraste mejorado se han consolidado como una herramienta esencial en el diagnóstico médico, permitiendo una representación más exacta de las estructuras internas del cuerpo y una detección superior de lesiones. La administración de agentes de contraste está específicamente diseñada para ``iluminar'' ciertas estructuras o áreas dentro del cuerpo humano, lo que conduce a una identificación más precisa de posibles problemas de salud.

La necesidad de técnicas efectivas de mejora del contraste en imágenes médicas es evidente debido a las dificultades que plantea el bajo contraste y la importancia de un diagnóstico certero. La investigación en esta área es, por lo tanto, de gran relevancia e impacto, ya que aborda un desafío fundamental que afecta directamente la capacidad de los profesionales médicos para realizar diagnósticos precisos.

En los últimos años, se ha observado una creciente aplicación de la inteligencia computacional, como los algoritmos genéticos, en el ámbito del mejoramiento de imágenes médicas. Si bien existen técnicas tradicionales de mejora del contraste, que incluyen la modificación del histograma y la corrección gamma, así como técnicas avanzadas como CLAHE (Ecualización de Histograma Adaptativa con Límite de Contraste), los algoritmos genéticos (AGs) ofrecen un enfoque alternativo para abordar la complejidad de este problema. Los AGs son técnicas de optimización metaheurísticas inspiradas en el proceso de selección natural, empleadas para encontrar soluciones óptimas a problemas complejos. Estos algoritmos resultan particularmente efectivos en problemas donde la función objetivo es discontinua, no diferenciable, estocástica o altamente no lineal, características que pueden presentarse en las tareas de procesamiento de imágenes. Los AGs operan sobre una población de soluciones candidatas que evolucionan a lo largo de generaciones mediante la aplicación de operadores de selección, cruce y mutación. La tendencia hacia el uso de la inteligencia computacional para la mejora de imágenes sugiere que los métodos tradicionales pueden tener limitaciones o que las técnicas más recientes ofrecen ventajas en escenarios específicos. Los algoritmos genéticos, con su capacidad para manejar problemas de optimización complejos, se presentan como una técnica avanzada viable.

El presente reporte se centra en el desarrollo de un marco teórico para un estudio que investiga el mejoramiento de contraste en imágenes médicas utilizando un algoritmo genético con transformación sigmoide. La investigación propuesta realiza una comparación entre la maximización de la entropía de la imagen y la maximización de la desviación estándar de la intensidad de los píxeles como funciones objetivo dentro del proceso de optimización. El estudio utilizará una función sigmoide para la transformación de la intensidad de la imagen, y el algoritmo genético optimizará los parámetros de esta transformación. El proceso de optimización estará guiado por dos funciones objetivo distintas: la maximización de la entropía de la imagen y la maximización de la desviación estándar de la intensidad de los píxeles. Este reporte proporcionará la base teórica para este enfoque y la comparación de estas dos funciones objetivo.

\chapter{Objetivos del Proyecto}

\section{Implementar un algoritmo genético robusto}
Desarrollar un Algoritmo Genético (AG) que utilice métodos avanzados de selección, cruzamiento y mutación para la optimización de funciones complejas.

\section{Minimización de funciones benchmark}
Aplicar el AG a las funciones Langermann y Drop-Wave, que presentan múltiples óptimos locales y retos propios en la búsqueda de la solución global.

\section{Visualización y análisis}
Generar salidas que permitan analizar la evolución del fitness a lo largo de las generaciones y visualizar la superficie de la función en 3D, facilitando la comprensión del comportamiento del algoritmo.

\chapter{Metodolog\'ia}

\section{Inicializaci\'on de la Poblaci\'on}
La poblaci\'on inicial se genera de forma uniforme a lo largo del espacio de b\'usqueda, definido por l\'imites inferiores y superiores para cada variable.

\subsection*{Objetivo}
Garantizar que la b\'usqueda comience explorando de manera equitativa todas las regiones posibles, evitando sesgos que puedan limitar la diversidad de soluciones iniciales.

\subsection*{Implementaci\'on}
Se utiliza la funci\'on \texttt{initialize\_population}, la cual emplea m\'etodos de generaci\'on aleatoria (por ejemplo, la funci\'on \texttt{np.random.uniform} de NumPy) para crear un conjunto de individuos.

\subsection*{Ventajas}
\begin{itemize}
    \item Permite cubrir todo el rango definido para cada variable.
    \item Aumenta la probabilidad de encontrar regiones prometedoras del espacio de soluciones desde el inicio.
\end{itemize}

\section{Evaluaci\'on de Fitness}

Cada individuo generado se eval\'ua mediante la funci\'on objetivo, la cual determina qu\'e tan buena es la soluci\'on propuesta.

\subsection*{Funciones Utilizadas}
\begin{itemize}
    \item \textbf{Entrop\'ia de Shannon}: Se utiliza para medir la cantidad de informaci\'on contenida en la imagen transformada. Un mayor valor de entrop\'ia indica una distribuci\'on de intensidades m\'as uniforme y, por lo tanto, un mejor contraste.
    \item \textbf{Desviaci\'on estandar}: Eval\'ua la dispersi\'on de los niveles de intensidad en la imagen. Una mayor desviaci\'on estándar sugiere una mayor separaci\'on entre tonos oscuros y claros, lo que contribuye a un mayor contraste.
\end{itemize}

\subsection*{Proceso}
Para cada individuo en la población, se aplican los valores de $\alpha$ y $\Delta$ $(x_1, x_2)$ a la transformada sigmoide, obteniendo una imagen modificada. Posteriormente se calculan las metricas de Desviaci\'on estandar y entrop\'ia  y se almacena dicho valor para posteriores comparaciones.

\subsection*{Importancia}
La evaluaci\'on correcta del fitness es crucial, ya que determina la selecci\'on de individuos y, por ende, el rumbo de la evoluci\'on poblacional.

\section{Selecci\'on por Torneo}

Para elegir los padres que generar\'an la siguiente generaci\'on se utiliza un m\'etodo de selecci\'on por torneo.

\subsection*{Mecanismo}
\begin{itemize}
    \item Se forman m\'ultiples grupos (torneos) de individuos seleccionados al azar.
    \item En cada grupo se compara el fitness de los participantes y se selecciona al individuo con el mejor desempe\~no.
\end{itemize}

\subsection*{Implementaci\'on Vectorizada}
La funci\'on \texttt{vectorized\_tournament\_selection} realiza este proceso de forma eficiente, aprovechando operaciones vectorizadas de NumPy.

\subsection*{Beneficios}
\begin{itemize}
    \item Favorece la selecci\'on de soluciones de alta calidad sin descartar por completo la diversidad poblacional.
    \item Permite controlar la presi\'on selectiva mediante el tama\~no del torneo.
\end{itemize}

\section{Cruzamiento con SBX}

El operador de cruzamiento se implementa mediante el m\'etodo SBX (Simulated Binary Crossover).

\subsection*{Proceso del SBX}
\begin{itemize}
    \item A partir de dos padres, se genera un n\'umero aleatorio $u$ y se calcula un par\'ametro $\beta$ que determina la dispersi\'on de los descendientes respecto a los padres.
    \item Se generan dos hijos combinando linealmente los valores de los padres.
\end{itemize}

\subsection*{Ajuste de L\'imites}
Se incorpora un mecanismo en \texttt{sbx\_crossover\_with\_boundaries} que garantiza que los hijos resultantes se mantengan dentro de los l\'imites predefinidos.

\subsection*{Ventajas}
\begin{itemize}
    \item Promueve la creaci\'on de soluciones intermedias que pueden explotar la informaci\'on gen\'etica de ambos padres.
    \item Ayuda a preservar la diversidad en la poblaci\'on.
\end{itemize}

\section{Mutaci\'on Polinomial}

Para introducir variabilidad y explorar nuevas regiones del espacio de b\'usqueda, se aplica la mutaci\'on polinomial.

\subsection*{Mecanismo de la Mutaci\'on}
\begin{itemize}
    \item Cada gen de un individuo tiene una probabilidad definida de sufrir una mutaci\'on.
    \item Se usa una distribuci\'on polinomial, controlada por el par\'ametro $\eta_{\text{mut}}$.
\end{itemize}

\subsection*{Consideraciones de L\'imites}
La mutaci\'on se aplica respetando los l\'imites definidos para cada variable mediante la funci\'on \texttt{polynomial\_mutation\_with\_boundaries}.

\subsection*{Beneficios}
\begin{itemize}
    \item Introduce peque\~nas variaciones que pueden conducir a la exploraci\'on de nuevas soluciones.
    \item Previene la convergencia prematura al mantener la diversidad gen\'etica.
\end{itemize}

\section{Elitismo y Ciclo Evolutivo}

El proceso evolutivo se estructura en ciclos o generaciones.

\subsection*{Elitismo}
\begin{itemize}
    \item Se retiene el mejor individuo de la generaci\'on actual y se garantiza su inclusi\'on en la siguiente generaci\'on.
    \item Esto asegura que la calidad de la soluci\'on nunca empeore a lo largo de las generaciones.
\end{itemize}

\subsection*{Ciclo Evolutivo}
\begin{itemize}
    \item Cada generaci\'on incluye la selecci\'on, el cruzamiento, la mutaci\'on y la incorporaci\'on del individuo de \'elite.
    \item La evoluci\'on se repite durante un n\'umero predefinido de generaciones.
\end{itemize}

\subsection*{Registro y An\'alisis}
\begin{itemize}
    \item Se almacena el historial del fitness y de las mejores soluciones.
    \item Esto facilita el an\'alisis del comportamiento del algoritmo y la generaci\'on de visualizaciones.
\end{itemize}

\chapter{Resultados y Discusión}

Durante la ejecución del algoritmo genético para la mejora de contraste en imágenes médicas, se realizaron múltiples corridas completas para cada función objetivo (Maximización de Entropía y Maximización de Desviación Estándar), lo que permitió evaluar la estabilidad y eficiencia del método. Los resultados se agruparon en resúmenes globales, donde se registraron indicadores clave, los cuales se encuentran dentro de las siguientes tablas:

\begin{table}[H]
\centering
\caption{Resumen Global de Corridas 01 - Imagen Médica 5R (Mano)}
\begin{tabular}{lcc}
\hline
\textbf{Indicador} & \textbf{Optimización de Entropía} & \textbf{Optimización de Desviación Estándar} \\
\hline
Mejor (Fitness) & -4.128421621722379 & -0.1953755630183538 \\
Peor (Fitness) & -4.128421621722379 & -0.1953755555673538 \\
Media & -4.128421621722379 & -0.1953755630183538 \\
Desv. Estándar & 0.0 & 7.51e-09 \\
\hline
\end{tabular}
\end{table}

\begin{table}[H]
    \centering
    \caption{Resumen Global de Corridas 02- Imagen Médica 5R (Mano)}
    \begin{tabular}{lcc}
    \hline
    \textbf{Indicador} & \textbf{Optimización de Desviación Estándar} \\
    \hline
    Mejor (Fitness) & 120.061897 \\
    Peor (Fitness) & 120.245979 \\
    Media & 120.14108900000001 \\
    Desv. Estándar & 0.07353604532132446 \\
    \hline
    \end{tabular}
\end{table}


\begin{table}[H]
    \centering
    \caption{Resumen Global de Corridas 01 - Imagen Médica 6R (Ojo)}
    \begin{tabular}{lcc}
    \hline
    \textbf{Indicador} & \textbf{Optimización de Entropía} & \textbf{Optimización de Desviación Estándar} \\
    \hline
    Mejor (Fitness) & -6.84574278505749 & -0.2244060136302129 \\
    Peor (Fitness) & -6.84574278505749 & -0.2244058045302129 \\
    Media & -6.84574278505749 & -0.2244060136302129 \\
    Desv. Estándar & 0.0 & 1.91e-05 \\
    \hline
    \end{tabular}
\end{table}

\begin{table}[H]
    \centering
    \caption{Resumen Global de Corridas 02 - Imagen Médica 6R (Ojo)}
    \begin{tabular}{lcc}
    \hline
    \textbf{Indicador} & \textbf{Optimización de Entropía} & \textbf{Optimización de Desviación Estándar} \\
    \hline
    Mejor (Fitness) & 4.7195 & 122.711517 \\
    Peor (Fitness) & 4.7211 & 122.773857 \\
    Media & 4.72044 & 122.7600892 \\
    Desv. Estándar & 0.0006066300355240354 & 0.027192434392680088 \\
    \hline
    \end{tabular}
\end{table}

Donde cada Indicador Representa lo siguiente:
\begin{itemize}
    \item \textbf{Mejor (Fitness):} Representa la solución con el valor de fitness mínimo obtenido en todas las corridas.
    \item \textbf{Peor (Fitness):} Indica la solución con el mayor valor de fitness, sirviendo como referencia de la variabilidad en la búsqueda.
    \item \textbf{Media:} Es el promedio de los valores de fitness de la mejor solución de cada corrida, ofreciendo una visión global del desempeño del algoritmo.
    \item \textbf{Desv. Estándar:} Mide la dispersión de los valores de fitness entre las corridas, reflejando la estabilidad y consistencia del proceso evolutivo.
\end{itemize}

\section{Análisis de los Resultados}

\subsection{Consistencia y Robustez}
Los resúmenes globales muestran que, a lo largo de las corridas, el algoritmo genético tiende a converger de manera consistente hacia soluciones de alta calidad. La extremadamente baja desviación estándar en los valores de fitness, especialmente para la optimización de entropía (0.0 en ambas imágenes), sugiere que el proceso evolutivo es excepcionalmente robusto y no depende en exceso de la aleatoriedad inherente a los operadores genéticos. Esto es crucial para problemas de optimización en imágenes médicas, ya que garantiza que la metodología aplicada es reproducible y confiable para aplicaciones clínicas.

\subsection{Comparación Entre Funciones Objetivo}
\subsubsection{Optimización de Entropía}
El resumen global para la optimización de entropía indica que el algoritmo fue capaz de identificar soluciones consistentes, con una desviación estándar de 0.0 en ambas imágenes. Para la Imagen 5R, se obtuvieron valores óptimos de $\alpha = 8.2005$ y $\delta \approx 0$, mientras que para la Imagen 6R, curiosamente, los valores óptimos fueron $\alpha \approx 0$ y $\delta \approx 0$. Esta diferencia significativa entre las dos imágenes sugiere que la optimización basada en entropía es altamente sensible a las características específicas de cada imagen médica.

\subsubsection{Optimización de Desviación Estándar}
Para la optimización de la desviación estándar, los resultados globales evidencian una convergencia hacia regiones con valores de fitness bajos, con una ligera variación entre corridas. Es notable que, para ambas imágenes, esta función objetivo condujo a valores de alpha cercanos al límite superior (aproximadamente 10), pero con diferentes valores de delta: cercano a 0 para la Imagen 5R y aproximadamente 0.41 para la Imagen 6R. Esto sugiere que la maximización de la desviación estándar tiende a favorecer transformaciones con pendientes pronunciadas, independientemente de la imagen específica.

\subsection{Parámetros Óptimos y Sensibilidad}
Los resultados muestran una clara diferencia en la sensibilidad de los parámetros:

\begin{itemize}
    \item \textbf{Parámetro Alpha:} Controla la pendiente de la transformación sigmoidea y, por tanto, la intensidad del contraste. La optimización de la desviación estándar consistentemente empuja este parámetro hacia su límite superior (10), mientras que la optimización de entropía muestra una mayor variabilidad dependiendo de la imagen.

    \item \textbf{Parámetro Delta:} Controla el desplazamiento del punto medio de la transformación, afectando al brillo general. Este parámetro muestra una mayor sensibilidad y especificidad según la imagen, especialmente para la optimización de la desviación estándar.
\end{itemize}

\subsection{Transformación Sigmoidea}
La mejora de contraste se logra mediante una transformación sigmoidea definida por la función:

\begin{equation}
sigmoid(x) = \frac{1}{1 + exp(-alpha \cdot (x - delta))}
\end{equation}

Donde los parámetros optimizados por el algoritmo genético son:
\begin{itemize}
    \item \textbf{alpha:} Controla la pendiente de la transformación (intensidad del contraste)
    \item \textbf{delta:} Controla el desplazamiento del punto medio (ajuste de brillo)
\end{itemize}

\section{Discusión de Resultados}

\begin{itemize}
    \item \textbf{Eficacia del Algoritmo:} Los indicadores globales extraídos de los CSV demuestran que el algoritmo genético es capaz de encontrar soluciones óptimas para la mejora de contraste en imágenes médicas, con una consistencia notable especialmente para la optimización de entropía.

    \item \textbf{Especificidad de la Imagen:} La significativa diferencia en los parámetros óptimos entre las imágenes 5R y 6R, particularmente para la optimización de entropía, indica que la mejora de contraste óptima es altamente dependiente de las características específicas de cada imagen. Esto respalda la necesidad de técnicas de mejora adaptativas e individualizadas en lugar de parámetros fijos.

    \item \textbf{Selección de Métricas:} La elección de la función objetivo impacta significativamente en los resultados de la mejora. La optimización de entropía parece estar mejor adaptada para preservar el contenido informativo, mientras que la optimización de la desviación estándar se centra en maximizar el contraste visual.

    \item \textbf{Equilibrio entre Exploración y Explotación:} La aplicación de operadores de selección por torneo, cruzamiento SBX y mutación polinomial, junto con una población de 100 individuos y 50 generaciones, ha demostrado ser efectiva para mantener un equilibrio entre la exploración del espacio de búsqueda y la explotación de las regiones prometedoras.

    \item \textbf{Potencial para Aplicaciones Clínicas:} La consistencia y robustez mostradas por los resultados sugieren que este enfoque puede ser viable para aplicaciones clínicas donde la reproducibilidad y la confiabilidad son cruciales para el diagnóstico médico.
\end{itemize}

\subsection{Visualización de Resultados}
Para complementar el análisis numérico, se generaron representaciones visuales de las imágenes originales y mejoradas, permitiendo una evaluación cualitativa de los efectos de la optimización.

\begin{figure}[H]
    \centering
    \includegraphics[width=\textwidth]{../Medical_Image-Contrast_Enhancement_using_Genetic_Algorithms/outputs/Original_Medica5R_Entropy/surface_3d_Original_Medica5R_Entropy.png}
    \caption{Comparación entre la imagen médica 5R original y su version mejorada mediante optimización de entropía}
    \label{fig:comparacion_5R_Entropia}
\end{figure}

\begin{figure}[H]
    \centering
    \includegraphics[width=\textwidth]{../Medical_Image-Contrast_Enhancement_using_Genetic_Algorithms/outputs/Original_Medica5R_Std_Deviation/surface_3d_Original_Medica5R_Std_Deviation.png}
    \caption{Comparación entre la imagen médica 5R original y su version mejoradas mediante optimización de desviación estándar}
    \label{fig:comparacion_5R_std}
\end{figure}

\begin{figure}[H]
    \centering
    \includegraphics[width=\textwidth]{../Medical_Image-Contrast_Enhancement_using_Genetic_Algorithms/outputs/Original_Medica6R_Entropy/surface_3d_Original_Medica6R_Entropy.png}
    \caption{Comparación entre la imagen médica 6R original y su versiones mejorada mediante optimización de entropía}
    \label{fig:comparacion_6R_Entropia}
\end{figure}

\begin{figure}[H]
    \centering
    \includegraphics[width=\textwidth]{../Medical_Image-Contrast_Enhancement_using_Genetic_Algorithms/outputs/Original_Medica6R_Std_Deviation/surface_3d_Original_Medica6R_Std_Deviation.png}
    \caption{Comparación entre la imagen médica 6R original y su versiones mejorada mediante optimización de entropía}
    \label{fig:comparacion_6R_Entropia}
\end{figure}
\newpage
\section{Pensamientos finales}

El algoritmo genético implementado ha demostrado ser eficaz para la optimización de parámetros de transformación sigmoidea en la mejora de contraste de imágenes médicas. Los resultados muestran que:

\begin{itemize}
    \item La optimización de entropía y la optimización de desviación estándar representan diferentes filosofías de mejora: la primera preserva el contenido informativo mientras que la segunda maximiza el contraste visible.

    \item Las diferencias significativas en los parámetros óptimos entre imágenes enfatizan la necesidad de técnicas de mejora adaptativas y específicas para cada imagen.

    \item La consistencia de convergencia a través de múltiples corridas indica la fiabilidad del algoritmo genético para encontrar soluciones óptimas.

    \item Un sistema ideal de mejora de imágenes médicas podría beneficiarse del uso de múltiples funciones objetivo o una métrica compuesta que equilibre la preservación de la información con la mejora del contraste.
\end{itemize}

Estos hallazgos sugieren que la aplicación de algoritmos genéticos para la optimización de parámetros de transformación en imágenes médicas proporciona un enfoque robusto y adaptable para mejorar la calidad de imagen, potencialmente conduciendo a diagnósticos más precisos en entornos clínicos.
\chapter{Implementación del Algoritmo Genético}
\label{chap:implementacion}

En este capítulo se detalla la implementación del algoritmo genético (GA) diseñado para optimizar el contraste de imágenes médicas, utilizando C++ y la biblioteca OpenCV. El GA ajusta los parámetros \(\alpha\) y \(\delta\) de una transformación sigmoidal, maximizando métricas como la entropía de Shannon o la desviación estándar. Se describen los fundamentos teóricos, la estructura del código principal y los operadores de cruce y mutación utilizados.

\section{Fundamento Teórico}
\label{sec:fundamento_teorico}

Un algoritmo genético es una técnica de optimización inspirada en los principios de la evolución biológica, como la selección natural, la reproducción y la mutación. Este método pertenece a la familia de los algoritmos evolutivos y se utiliza para resolver problemas complejos de optimización donde los métodos tradicionales pueden ser insuficientes, como en la mejora del contraste de imágenes médicas.

El GA opera sobre una población de posibles soluciones, denominadas individuos. Cada individuo representa una combinación de parámetros que se desea optimizar. En este caso, los parámetros son \(\alpha\) y \(\delta\), que controlan una transformación sigmoidal aplicada a las intensidades de los píxeles de una imagen. El objetivo es encontrar los valores óptimos de estos parámetros que maximicen una métrica específica de calidad de la imagen, como la entropía de Shannon o la desviación estándar.

El proceso evolutivo se desarrolla en iteraciones conocidas como generaciones, y en cada una de ellas se aplican los siguientes pasos:

\begin{itemize}
    \item \textbf{Inicialización:} Se genera una población inicial de individuos con valores aleatorios dentro de rangos predefinidos (por ejemplo, \(\alpha \in [0, 20]\) y \(\delta \in [0, 1]\)).
    \item \textbf{Evaluación de la aptitud:} Se calcula una función de aptitud para cada individuo, que mide qué tan buena es la solución que representa. En este caso, la aptitud puede ser la entropía o la desviación estándar de la imagen transformada.
    \item \textbf{Selección:} Se eligen los individuos más aptos (con mayor valor de aptitud) para que actúen como "padres" de la siguiente generación.
    \item \textbf{Cruce:} Se combinan pares de individuos seleccionados para producir "descendientes", intercambiando información genética (valores de \(\alpha\) y \(\delta\)).
    \item \textbf{Mutación:} Se introducen pequeñas variaciones aleatorias en algunos individuos para mantener la diversidad en la población y evitar converger prematuramente a una solución subóptima.
    \item \textbf{Reemplazo:} La nueva población reemplaza a la anterior, preservando a veces al mejor individuo (elitismo), y el ciclo se repite hasta alcanzar un número fijo de generaciones o un criterio de parada.
\end{itemize}

\subsection{Transformación Sigmoidal}
El núcleo del problema radica en ajustar los parámetros de una transformación sigmoidal que mejora el contraste de la imagen. La ecuación de esta transformación es:

\begin{equation}
    I^\prime(x, y) = \frac{1}{1 + e^{-\alpha (I(x, y) - \delta)}}
\label{eq:sigmoid_transform}
\end{equation}

donde:

\begin{itemize}
    \item \(I(x, y)\) es la intensidad original del píxel en la posición \((x, y)\) de la imagen en escala de grises (normalizada entre 0 y 1).
    \item \(I^\prime(x, y)\) es la intensidad transformada del píxel.
    \item \(\alpha\) controla la pendiente de la curva sigmoidal, determinando la brusquedad del cambio de intensidad.
    \item \(\delta\) define el punto de inflexión de la curva, desplazando el rango de intensidades afectado.
\end{itemize}

Esta función transforma las intensidades de la imagen de manera no lineal, amplificando el contraste al mapear los valores originales a un rango que resalta diferencias entre regiones claras y oscuras.

\subsection{Función de Aptitud}
La función de aptitud es el criterio que guía la optimización. En este caso, se utilizan dos métricas comunes para evaluar el contraste de la imagen transformada:

\textbf{Entropía de Shannon:}
\begin{equation}
    H = -\sum_{i=0}^{255} p_i \log_2(p_i)
\end{equation}
donde \(p_i\) es la probabilidad de la intensidad \(i\). Una mayor entropía indica una distribución más uniforme de intensidades, lo que suele corresponder a un mejor contraste.

\textbf{Desviación Estándar:}
\begin{equation}
    \sigma = \sqrt{\frac{1}{N} \sum_{i=1}^{N} (x_i - \mu)^2}
\end{equation}
donde \(x_i\) son las intensidades de los píxeles, \(\mu\) es la media y \(N\) es el número total de píxeles. Una desviación estándar más alta implica mayor variabilidad en las intensidades, lo que también se asocia con un contraste mejorado.

El GA busca maximizar estas métricas, aunque en la implementación se minimiza el valor negativo de la entropía o la desviación estándar para adaptarse a la convención de minimización en muchos algoritmos.

\section{Estructura del Código Principal}\label{sec:estructura_codigo}

El núcleo del algoritmo está implementado en dos archivos: \texttt{genetic\_algo\_img.h} y \texttt{genetic\_algo\_img.cpp}. El primero define la clase \texttt{GeneticAlgorithmReal}, mientras que el segundo contiene la lógica de sus métodos.

\subsection{Definición de la Clase}\label{subsec:definicion_clase}

En \texttt{genetic\_algo\_img.h}, se declara la clase con los elementos esenciales para el GA:

\begin{lstlisting}[style=cppstyle, caption={Definición de la clase en genetic\_algo\_img.h}, label={lst:ga_header}]
class GeneticAlgorithmReal {
public:
    enum class FitnessMetric { ENTROPY, STDDEV };
private:
    int population_size;
    int num_generations;
    std::vector<double> lower_bound;
    std::vector<double> upper_bound;
    cv::Mat gray_img;
    std::vector<std::vector<double>> population;
    std::vector<double> fitness_values;
    FitnessMetric metric;

    void initialize_population();
    void evaluate_fitness();
    void crossover();
    void mutate();
public:
    GeneticAlgorithmReal(int pop_size, int num_genes, int num_generations,
                         double crossover_prob, double mutation_prob,
                         std::vector<double> lower_bound,
                         std::vector<double> upper_bound,
                         const cv::Mat& gray_img,
                         FitnessMetric metric);
    void operator()();
    std::pair<std::vector<double>, double> getBestSolution();
};
\end{lstlisting}

La enumeración \texttt{FitnessMetric} permite alternar entre entropía y desviación estándar como criterios de aptitud. Los miembros privados incluyen parámetros del GA (tamaño de población, generaciones, límites) y datos como la imagen de entrada y la población de individuos.

\subsection{Inicialización de la Población}\label{subsec:inicializacion}

El método \texttt{initialize\_population()} genera una población inicial de individuos aleatorios dentro de los límites establecidos:

\begin{lstlisting}[style=cppstyle, caption={Inicialización de la población}, label={lst:init_pop}]
void GeneticAlgorithmReal::initialize_population() {
    for (auto& individual : population) {
        for (int i = 0; i < num_genes; ++i) {
            individual[i] = lower_bound[i] + dis(gen) * (upper_bound[i] - lower_bound[i]);
            individual[i] = std::max(lower_bound[i], std::min(upper_bound[i], individual[i]));
        }
    }
}
\end{lstlisting}

La función \texttt{dis(gen)} genera valores aleatorios uniformes, y las funciones \texttt{std::max} y \texttt{std::min} aseguran que \(\alpha\) y \(\delta\) permanezcan dentro de \([0, 20]\) y \([0, 1]\), respectivamente.

\subsection{Evaluación de la Aptitud}
\label{subsec:evaluacion_aptitud}

El método \texttt{evaluate\_fitness()} calcula la aptitud de cada individuo aplicando la transformación sigmoidal:

\begin{lstlisting}[style=cppstyle, caption={Evaluación de la aptitud}, label={lst:eval_fitness}]
void GeneticAlgorithmReal::evaluate_fitness() {
    for (int i = 0; i < population_size; ++i) {
        double alpha = population[i][0];
        double delta = population[i][1];
        cv::Mat transformed = sigmoid_transform(gray_img, alpha, delta);
        cv::normalize(transformed, transformed, 0, 1, cv::NORM_MINMAX);
        transformed.convertTo(transformed, CV_8U, 255.0);
        if (metric == FitnessMetric::ENTROPY) {
            double entropy_value = entropy(transformed);
            fitness_values[i] = -entropy_value; // Minimize negative entropy (maximize entropy)
        } else { // STDDEV
            double stddev_value = standard_deviation(transformed);
            fitness_values[i] = -stddev_value; // Minimize negative stddev (maximize stddev)
        }
    }
}
\end{lstlisting}

La imagen transformada se normaliza y convierte a formato de 8 bits. La aptitud se define como el negativo de la métrica seleccionada, permitiendo al GA maximizarla mediante un proceso de minimización.

\section{Operadores Genéticos}
\label{sec:operadores}

El GA utiliza dos operadores clave: el cruce SBX y la mutación polinomial, implementados en \texttt{sbx.cpp} y \texttt{polynomial\_mut.cpp}, respectivamente.

\subsection{Cruce SBX}
\label{subsec:cruce_sbx}

El operador Simulated Binary Crossover (SBX) genera dos descendientes a partir de dos padres, simulando un cruce binario en variables reales:

\begin{lstlisting}[style=cppstyle, caption={Cruce SBX}, label={lst:sbx}]
std::pair<std::vector<double>, std::vector<double>> sbx(std::vector<double> f1, std::vector<double> f2, double r,
    const std::vector<double>& lower_bound, const std::vector<double>& upper_bound, double nc) {
    std::vector<double> s1(f1.size());
    std::vector<double> s2(f2.size());
    for(size_t i = 0; i < f1.size(); i++) {
        if (std::abs(f1[i] - f2[i]) < 1e-14) {
            s1[i] = f1[i];
            s2[i] = f2[i];
            continue;
        }
        double beta = 1 + (2.0 * std::min(f1[i] - lower_bound[i], upper_bound[i] - f2[i]) /
                           std::abs(f1[i] - f2[i]));
        double alpha = 2 - std::pow(std::abs(beta), -(nc + 1.0));
        double u = r;
        double beta_c = (u <= 1.0 / alpha) ?
                        std::pow((u * alpha), (1.0 / (nc + 1))) :
                        std::pow((1.0 / (2.0 - u * alpha)), (1.0 / (nc + 1)));
        s1[i] = 0.5 * (f1[i] + f2[i] - beta_c * std::abs(f2[i] - f1[i]));
        s2[i] = 0.5 * (f1[i] + f2[i] + beta_c * std::abs(f2[i] - f1[i]));
    }
    return {s1, s2};
}
\end{lstlisting}

El parámetro \(\beta_c\) controla la dispersión de los descendientes (\(s1\), \(s2\)) respecto a los padres (\(f1\), \(f2\)), basado en una distribución simulada binaria ajustada por el índice \(nc\).

\subsection{Mutación Polinomial}
\label{subsec:mutacion_polinomial}

La mutación polinomial introduce variaciones controladas en los parámetros:

\begin{lstlisting}[style=cppstyle, caption={Mutación polinomial}, label={lst:poly_mut}]
std::vector<double> polynomial_mutation(std::vector<double>& x, double r, int nm) {
    std::vector<double> mutated_x(x.size());
    for (int i = 0; i < x.size(); i++) {
        double upper_bound = x[i] > 0 ? std::floor(x[i]) : std::ceil(x[i]);
        double lower_bound = x[i] < 0 ? std::floor(x[i]) : std::ceil(x[i]);
        double delta = std::min(upper_bound - x[i], x[i] - lower_bound) / (upper_bound - lower_bound);
        double deltaq = 0.0;
        if (r <= 0.5) {
            deltaq = x[i] + std::pow(2 * delta + (1 - 2 * delta) * (1 - r), nm + 1);
        } else {
            deltaq = x[i] - std::pow(2 * (1 - delta) + 2 * (delta - 0.5) * (1 - r), nm + 1);
        }
        mutated_x[i] = deltaq * (upper_bound - lower_bound);
    }
    return mutated_x;
}
\end{lstlisting}

El valor \(\delta_q\) representa la perturbación, cuya magnitud depende del índice de mutación \(nm\), permitiendo ajustes finos en los parámetros.

\section{Ejecución del Algoritmo}
\label{sec:ejecucion}

El método \texttt{operator()()} ejecuta el ciclo principal del GA:

\begin{lstlisting}[style=cppstyle, caption={Ciclo principal del algoritmo}, label={lst:ga_run}]
void GeneticAlgorithmReal::operator()() {
    for (int generation = 0; generation < num_generations; ++generation) {
        std::cout << "Generation " << generation + 1 << std::endl;
        output_file << "Generation: " << generation + 1 << " | ";
        evaluate_fitness();
        select_best_individual();
        crossover();
        mutate();
        evaluate_fitness();
        apply_elitism();
    }
}
\end{lstlisting}

Este ciclo itera sobre las generaciones, evaluando la aptitud, seleccionando individuos, aplicando los operadores SBX y mutación polinomial,
y usando elitismo a fines de preservar la mejor solución.

\chapter{Conclusiones}
El proyecto demuestra la efectividad de los algoritmos genéticos en la optimización de funciones complejas. La implementación de técnicas avanzadas, como el cruzamiento SBX y la mutación polinomial, combinada con una estrategia de selección por torneo, ha permitido explorar de manera eficiente el espacio de soluciones y mejorar progresivamente el fitness de la población. Las herramientas de visualización y el almacenamiento de resultados facilitan el análisis del comportamiento del algoritmo y ofrecen una base sólida para futuras mejoras o aplicaciones a problemas más complejos.

Además, el algoritmo genético desarrollado constituye un marco sólido y escalable para la optimización de funciones complejas. Los resultados obtenidos respaldan la viabilidad del enfoque y abren la puerta a futuras investigaciones, ya sea para afinar los parámetros del algoritmo o para extender su aplicación a problemas con mayores dimensiones o características más complejas. Esta base permite, además, la incorporación de mejoras y la adaptación del método a diferentes contextos, consolidando su utilidad en el ámbito de la optimización computacional.
\appendix
\chapter{GitHub}
Se anexa un enlace al repositorio donde se encuentran los codigos y una README que expande la información aqui mostrada: \href{https://github.com/CARBAJE/Practica01-TSAB}{\textit{Click aqui para abrir el enlace al repositorio}}
%funcion\chapter{Scripts}

\section{Archivo \texttt{main.py}}
\begin{lstlisting}[
  caption={Implementación de \texttt{main.py}},
  style=pythonstyle,
  basicstyle=\ttfamily\footnotesize
]
import os
import pandas as pd
from AG_confs import *

from AG import genetic_algorithm
from libs.plot import *

def main():
    # Crear carpetas de salida generales
    os.makedirs("outputs", exist_ok=True)
    
    for func_key, func_data in FUNCTIONS.items():
        f_obj = func_data["func"]
        lb = func_data["lb"]
        ub = func_data["ub"]
        func_name = func_data["name"]
        num_runs = func_data["num_runs"]
        
        # Carpetas especificas de cada funcion
        func_folder = f"outputs/{func_key}"
        os.makedirs(func_folder, exist_ok=True)
        hist_folder = os.path.join(func_folder, "historiales")
        res_folder = os.path.join(func_folder, "resumenes")
        os.makedirs(hist_folder, exist_ok=True)
        os.makedirs(res_folder, exist_ok=True)
        
        print(f"\n==============================================")
        print(f"  FUNCIoN: {func_name}")
        print(f"==============================================")
        
        all_runs_history = []
        best_solutions_all_runs = []  # Guardaremos los mejores individuos (x1, x2) de cada corrida
        best_values_across_runs = []  # Guardaremos el best_val (fitness) de cada corrida
        
        for run in range(num_runs):
            print(f"\nEjecucion {run+1}/{num_runs}")
            
            (best_sol, best_val,
             worst_sol, worst_val,
             avg_sol,  avg_val,
             std_val,
             best_fitness_history,
             best_x1_history,
             best_x2_history,
             population_final,
             fitness_final,
             best_solutions_over_time) = genetic_algorithm(
                 f_obj, lb, ub,
                 pop_size=POP_SIZE,
                 num_generations=NUM_GENERATIONS,
                 tournament_size=TOURNAMENT_SIZE,
                 crossover_prob=CROSSOVER_PROB,
                 eta_c=ETA_C,
                 mutation_prob=MUTATION_PROB,
                 eta_mut=ETA_MUT
             )
            
            # 1) Guardar historial
            df_historial = pd.DataFrame({
                "Generacion": np.arange(1, NUM_GENERATIONS + 1),
                "Mejor x1": best_x1_history,
                "Mejor x2": best_x2_history,
                "Mejor Fitness": best_fitness_history
            })
            historial_filename = os.path.join(hist_folder, f"historial_run_{run+1}.csv")
            df_historial.to_csv(historial_filename, index=False)
            
            # 2) Guardar resumen de la corrida
            data_resumen = [
                ["Mejor", best_sol[0], best_sol[1], best_val],
                ["Media", avg_sol[0], avg_sol[1], avg_val],
                ["Peor", worst_sol[0], worst_sol[1], worst_val],
                ["Desv. estandar", np.nan, np.nan, std_val]
            ]
            df_resumen = pd.DataFrame(data_resumen, columns=["Indicador", "x1", "x2", "Fitness"])
            resumen_filename = os.path.join(res_folder, f"resumen_run_{run+1}.csv")
            df_resumen.to_csv(resumen_filename, index=False)
            
            print(df_resumen.to_string(index=False))
            
            all_runs_history.append(best_fitness_history)
            best_solutions_all_runs.append(best_sol)
            best_values_across_runs.append(best_val)
        
        # ===========================================
        #       RESUMEN GLOBAL DE LAS CORRIDAS
        # ===========================================
        best_values_arr = np.array(best_values_across_runs)
        solutions_arr = np.array(best_solutions_all_runs)  # Cada fila: [x1, x2] del mejor individuo de cada corrida
        
        # Para el "Mejor" y "Peor", buscamos el indice de la corrida con minimo y maximo fitness
        min_index = np.argmin(best_values_arr)
        max_index = np.argmax(best_values_arr)
        
        data_global = [
            ["Mejor (Fitness)", solutions_arr[min_index, 0], solutions_arr[min_index, 1], best_values_arr[min_index]],
            ["Peor (Fitness)", solutions_arr[max_index, 0], solutions_arr[max_index, 1], best_values_arr[max_index]],
            ["Media", np.mean(solutions_arr[:, 0]), np.mean(solutions_arr[:, 1]), np.mean(best_values_arr)],
            ["Desv. Estandar", np.std(solutions_arr[:, 0]), np.std(solutions_arr[:, 1]), np.std(best_values_arr)]
        ]
        df_global = pd.DataFrame(data_global, columns=["Indicador", "x1", "x2", "Fitness"])
        
        global_filename = os.path.join(res_folder, "resumen_global_corridas.csv")
        df_global.to_csv(global_filename, index=False)
        
        # (Opcional) Graficar evolucion del fitness de todas las corridas
        plot_evolucion_fitness(all_runs_history, func_key, func_name)
        
        # (Opcional) Graficar superficie 3D si la funcion es de 2 variables
        if len(lb) == 2:
            plot_surface_3d(f_obj, lb, ub, best_solutions_all_runs, func_key, func_name)

if __name__ == "__main__":
    main()

\end{lstlisting}

\section{Archivo \texttt{AG\_confs.py}}
\begin{lstlisting}[
  caption={Implementación de \texttt{AG\_confs.py}},
  style=pythonstyle,
  basicstyle=\ttfamily\footnotesize
]
import numpy as np
from libs.functions import langermann, drop_wave
# ---------------------------
# Parametros del algoritmo
# ---------------------------
POP_SIZE = 100            # Numero de individuos en la poblacion
NUM_GENERATIONS = 200     # Numero de generaciones
NUM_RUNS = 10             # Numero de ejecuciones completas (ciclos)

# Parametros de la funcion de Langermann
a = np.array([3, 5, 2, 1, 7])
b = np.array([5, 2, 1, 4, 9])
c = np.array([1, 2, 5, 2, 3])

# Parametros del torneo
TOURNAMENT_SIZE = 3  # Numero de individuos participantes en cada torneo

# Parametros del cruzamiento SBX
CROSSOVER_PROB = 0.9  # Probabilidad de aplicar cruzamiento
ETA_C = 15            # indice de distribucion para SBX

# Parametros de la mutacion polinomial
MUTATION_PROB = 10.0 / 2  # Probabilidad de mutar cada gen
ETA_MUT = 20              # indice de distribucion para mutacion polinomial

best_solutions_list = [] 
all_runs_history = []  # Para graficar luego

FUNCTIONS = {
    "langermann": {
        "func": langermann,
        "lb": np.array([0, 0]),
        "ub": np.array([10, 10]),
        "name": "Langermann",
        "num_runs": NUM_RUNS
    },
    "drop_wave": {
        "func": drop_wave,
        "lb": np.array([-5.12, -5.12]),
        "ub": np.array([5.12, 5.12]),
        "name": "Drop-Wave",
        "num_runs": NUM_RUNS
    }
}
\end{lstlisting}

\section{Archivo \texttt{AG.py}}
\begin{lstlisting}[
  caption={Implementación de \texttt{AG.py}},
  style=pythonstyle,
  basicstyle=\ttfamily\footnotesize
]
from AG_confs import *
from libs.selection import vectorized_tournament_selection
from libs.crossover import sbx_crossover_with_boundaries
from libs.mutation import polynomial_mutation_with_boundaries
from libs.auxiliaries_functions import initialize_population


# ---------------------------
# Funcion principal del GA
# ---------------------------
def genetic_algorithm(objective_func, lower_bound, upper_bound,
                      pop_size=POP_SIZE, num_generations=NUM_GENERATIONS,
                      tournament_size=TOURNAMENT_SIZE,
                      crossover_prob=CROSSOVER_PROB, eta_c=ETA_C,
                      mutation_prob=MUTATION_PROB, eta_mut=ETA_MUT):
    """
    Ejecuta el GA para la funcion objetivo dada y retorna:
      - best_solution, best_value
      - worst_solution, worst_value
      - avg_solution, avg_value
      - std_value (fitness)
      - best_fitness_history, best_x1_history, best_x2_history
      - population (final), fitness (final)
      - best_solutions_over_time (para animaciones)
    """
    num_variables = len(lower_bound)
    
    # 1) Inicializar poblacion
    population = initialize_population(pop_size, num_variables, lower_bound, upper_bound)
    fitness = np.array([objective_func(ind) for ind in population])
    
    best_fitness_history = []
    best_x1_history = []
    best_x2_history = []
    
    # Para animacion: almacenamos el mejor (x1, x2) en cada generacion
    best_solutions_over_time = np.zeros((num_generations, num_variables))
    
    for gen in range(num_generations):
        # Elitismo: guardar el mejor de la generacion actual
        best_index = np.argmin(fitness)
        best_fitness = fitness[best_index]
        elite = population[best_index].copy()
        
        best_fitness_history.append(best_fitness)
        best_x1_history.append(elite[0])
        best_x2_history.append(elite[1])
        best_solutions_over_time[gen, :] = elite
        
        new_population = []
        
        # Numero de padres necesarios (2 por cada par a generar)
        num_parents_needed = 2 * (pop_size - 1)
        winners, _ = vectorized_tournament_selection(fitness, num_parents_needed,
                                                     tournament_size, len(population),
                                                     unique_in_column=True, unique_in_row=False)
        
        # Generar un valor global para el crossover y otro para la mutacion (para toda la generacion)
        global_u = np.random.rand()
        global_r = np.random.rand()
        
        # Generar nueva poblacion
        for i in range(0, len(winners), 2):
            parent1 = population[winners[i]].copy()
            if i + 1 < len(winners):
                parent2 = population[winners[i+1]].copy()
            else:
                parent2 = parent1.copy()
            
            # Cruzamiento SBX usando el mismo u para todas las variables del cruce
            child1, child2 = sbx_crossover_with_boundaries(
                parent1, parent2, lower_bound, upper_bound,
                eta_c, crossover_prob, use_global_u=True, global_u=global_u
            )
            # Mutacion polinomial usando el mismo r para todas las variables del individuo
            child1 = polynomial_mutation_with_boundaries(
                child1, lower_bound, upper_bound,
                mutation_prob, eta_mut, use_global_r=True, global_r=global_r
            )
            child2 = polynomial_mutation_with_boundaries(
                child2, lower_bound, upper_bound,
                mutation_prob, eta_mut, use_global_r=True, global_r=global_r
            )
            
            new_population.append(child1)
            if len(new_population) < pop_size - 1:
                new_population.append(child2)
        
        # Convertir a array y evaluar el fitness de la nueva poblacion
        new_population = np.array(new_population)
        new_fitness = np.array([objective_func(ind) for ind in new_population])
        
        # Incorporar el individuo elite (elitismo)
        new_population = np.vstack([new_population, elite])
        new_fitness = np.append(new_fitness, best_fitness)
        
        # Actualizar la poblacion y su fitness para la siguiente generacion
        population = new_population.copy()
        fitness = new_fitness.copy()
    
    # Calcular estadisticas finales
    best_index = np.argmin(fitness)
    worst_index = np.argmax(fitness)
    best_solution = population[best_index]
    best_value = fitness[best_index]
    worst_solution = population[worst_index]
    worst_value = fitness[worst_index]
    avg_solution = np.mean(population, axis=0)
    avg_value = np.mean(fitness)
    std_value = np.std(fitness)
    
    return (best_solution, best_value,
            worst_solution, worst_value,
            avg_solution, avg_value,
            std_value,
            best_fitness_history,
            best_x1_history,
            best_x2_history,
            population,
            fitness,
            best_solutions_over_time)
\end{lstlisting}

\section{Archivo \texttt{selection.py}}
\begin{lstlisting}[
  caption={Implementación de \texttt{selection.py}},
  style=pythonstyle,
  basicstyle=\ttfamily\footnotesize
]
import numpy as np

def vectorized_tournament_selection(fitness, num_tournaments, tournament_size, pop_size,
                                    unique_in_column=True, unique_in_row=False):
    """
    Genera una matriz de torneos de forma vectorizada y retorna, para cada torneo,
    el indice del individuo ganador (el de menor fitness).
    
    Args:
      - fitness: array con los fitness de la poblacion (longitud = pop_size).
      - num_tournaments: numero de torneos a realizar (por ejemplo, el numero total
                         de selecciones de padres requeridas en la generacion).
      - tournament_size: numero de individuos que participan en cada torneo.
      - pop_size: tamano de la poblacion.
      - unique_in_column: si True, para cada posicion (columna) se eligen candidatos sin
                          repeticion entre torneos.
      - unique_in_row: si True, en cada torneo (fila) los candidatos seran unicos.
                    (Por defecto se permite repetir en la fila).
    
    Returns:
      - winners: array de indices ganadores (uno por torneo).
      - tournament_matrix: la matriz de candidatos (de tamano [num_tournaments x tournament_size]).
    """
    if unique_in_row:
        # Para cada torneo (fila), muestreamos sin reemplazo (cada fila es unica)
        tournament_matrix = np.array([np.random.choice(pop_size, size=tournament_size, replace=False)
                                      for _ in range(num_tournaments)])
    else:
        # Permitir repeticion en la fila, pero controlar la no repeticion en cada columna
        if unique_in_column:
            # Para cada columna, se genera una permutacion de los indices (o se usan numeros aleatorios sin repeticion)
            # Siempre que num_tournaments <= pop_size.
            if num_tournaments > pop_size:
                # Si se requieren mas torneos que individuos, se hace sin la restriccion por columna.
                tournament_matrix = np.random.randint(0, pop_size, size=(num_tournaments, tournament_size))
            else:
                cols = []
                for j in range(tournament_size):
                    # Para la columna j, se toman num_tournaments indices sin repeticion
                    perm = np.random.permutation(pop_size)
                    cols.append(perm[:num_tournaments])
                tournament_matrix = np.column_stack(cols)
        else:
            # Sin restricciones, se muestrea con reemplazo para cada candidato.
            tournament_matrix = np.random.randint(0, pop_size, size=(num_tournaments, tournament_size))
    
    # Para cada torneo (fila de la matriz), se selecciona el candidato con el menor fitness.
    winners = []
    for row in tournament_matrix:
        row_fitness = fitness[row]
        winner_index = row[np.argmin(row_fitness)]
        winners.append(winner_index)
    winners = np.array(winners)
    return winners, tournament_matrix
\end{lstlisting}

\section{Archivo \texttt{crossover.py}}
\begin{lstlisting}[
  caption={Implementación de \texttt{crossover.py}},
  style=pythonstyle,
  basicstyle=\ttfamily\footnotesize
]
import numpy as np

def sbx_crossover(parent1, parent2, lower_bound, upper_bound, eta, crossover_prob):
    """Realiza el cruzamiento SBX para dos padres y devuelve dos hijos."""
    child1 = np.empty_like(parent1)
    child2 = np.empty_like(parent2)
    
    if np.random.rand() <= crossover_prob:
        for i in range(len(parent1)):
            u = np.random.rand()
            if u <= 0.5:
                beta = (2*u)**(1/(eta+1))
            else:
                beta = (1/(2*(1-u)))**(1/(eta+1))
            
            # Genera los dos hijos
            child1[i] = 0.5*((1+beta)*parent1[i] + (1-beta)*parent2[i])
            child2[i] = 0.5*((1-beta)*parent1[i] + (1+beta)*parent2[i])
            
            # Asegurar que los hijos esten dentro de los limites
            child1[i] = np.clip(child1[i], lower_bound[i], upper_bound[i])
            child2[i] = np.clip(child2[i], lower_bound[i], upper_bound[i])
    else:
        child1 = parent1.copy()
        child2 = parent2.copy()
    
    return child1, child2

def sbx_crossover_with_boundaries(parent1, parent2, lower_bound, upper_bound,
                                  eta, crossover_prob, use_global_u=False, global_u=None):
    """
    Realiza el cruzamiento SBX con limites, usando formulas que ajustan beta en funcion
    de la cercania a las fronteras. Permite usar un unico 'u' global para todos los individuos 
    de la generacion o, de forma estandar, un 'u' distinto por cada gen.
    
    Args:
      - parent1, parent2: arrays con los padres.
      - lower_bound, upper_bound: arrays con los limites inferiores y superiores.
      - eta: indice de distribucion para SBX.
      - crossover_prob: probabilidad de aplicar el cruce.
      - use_global_u: si es True se utilizara el mismo valor de 'u' para todas las variables.
      - global_u: valor de 'u' que se aplicara globalmente (si se proporciona).
      
    Returns:
      - child1, child2: arrays con los hijos resultantes.
    """
    parent1 = np.asarray(parent1)
    parent2 = np.asarray(parent2)
    child1 = np.empty_like(parent1)
    child2 = np.empty_like(parent2)
    
    # Si no se realiza el crossover, retornamos copias de los padres.
    if np.random.rand() > crossover_prob:
        return parent1.copy(), parent2.copy()
    
    # Si se quiere usar un 'u' global y no se ha pasado, se genera uno.
    if use_global_u:
        if global_u is None:
            global_u = np.random.rand()
    
    for i in range(len(parent1)):
        x1 = parent1[i]
        x2 = parent2[i]
        lb = lower_bound[i]
        ub = upper_bound[i]
        
        # Aseguramos que x1 sea menor o igual que x2
        if x1 > x2:
            x1, x2 = x2, x1
        
        dist = x2 - x1
        if dist < 1e-14:
            child1[i] = x1
            child2[i] = x2
            continue
        
        # Calcular la minima distancia a las fronteras
        min_val = min(x1 - lb, ub - x2)
        if min_val < 0:
            min_val = 0
        
        beta = 1.0 + (2.0 * min_val / dist)
        alpha = 2.0 - beta**(-(eta+1))
        
        # Si se usa u global, se usa el mismo valor para cada variable
        if use_global_u:
            u = global_u
        else:
            u = np.random.rand()
        
        if u <= (1.0 / alpha):
            betaq = (alpha * u)**(1.0/(eta+1))
        else:
            betaq = (1.0 / (2.0 - alpha*u))**(1.0/(eta+1))
        
        # Calcular los hijos
        c1 = 0.5 * ((x1 + x2) - betaq * (x2 - x1))
        c2 = 0.5 * ((x1 + x2) + betaq * (x2 - x1))
        
        # Ajustar a los limites
        child1[i] = np.clip(c1, lb, ub)
        child2[i] = np.clip(c2, lb, ub)
    
    return child1, child2
\end{lstlisting}

\section{Archivo \texttt{mutation.py}}
\begin{lstlisting}[
  caption={Implementación de \texttt{mutation.py}},
  style=pythonstyle,
  basicstyle=\ttfamily\footnotesize
]
import numpy as np

def polynomial_mutation(child, lower_bound, upper_bound, mutation_prob, eta_mut):
    """Aplica mutacion polinomial a un hijo."""
    mutant = child.copy()
    for i in range(len(child)):
        if np.random.rand() < mutation_prob:
            r = np.random.rand()
            diff = upper_bound[i] - lower_bound[i]
            if r < 0.5:
                delta = (2*r)**(1/(eta_mut+1)) - 1
            else:
                delta = 1 - (2*(1-r))**(1/(eta_mut+1))
            mutant[i] = child[i] + delta * diff
            mutant[i] = np.clip(mutant[i], lower_bound[i], upper_bound[i])
    return mutant


def polynomial_mutation_with_boundaries(child, lower_bound, upper_bound,
                                        mutation_prob, eta_mut,
                                        use_global_r=False, global_r=None):
    """
    Aplica mutacion polinomial (con limites) a un vector 'child'.
    Puede usar un unico 'r' global para todas las variables (si use_global_r=True)
    o generar un 'r' distinto para cada variable.
    
    Args:
      - child : array-like
          Cromosoma (vector de decision) a mutar.
      - lower_bound, upper_bound : array-like
          Limites inferiores y superiores para cada variable.
      - mutation_prob : float
          Probabilidad de mutacion (en [0,1]) para cada variable.
      - eta_mut : float
          indice de distribucion para la mutacion.
      - use_global_r : bool
          Si True, se utiliza un unico valor 'r' para todas las variables.
      - global_r : float, opcional
          Valor de 'r' global a usar; si no se proporciona, se genera uno.
    
    Retutrns:
      - mutant : np.ndarray
          Nuevo vector mutado (manteniendo la dimension de 'child').
    """
    mutant = np.array(child, copy=True, dtype=float)
    num_vars = len(child)
    
    # Si se desea usar un 'r' global y no se ha proporcionado, se genera uno una sola vez.
    if use_global_r:
        if global_r is None:
            global_r = np.random.rand()
    
    for i in range(num_vars):
        # Decidir si mutar esta variable
        if np.random.rand() < mutation_prob:
            x = mutant[i]
            xl = lower_bound[i]
            xu = upper_bound[i]
            
            # Evitar division por cero si los limites son casi iguales
            if abs(xu - xl) < 1e-14:
                continue

            # d = distancia normalizada al limite mas cercano
            d = min(xu - x, x - xl) / (xu - xl)
            
            # Elegir r: global o individual para cada variable
            if use_global_r:
                r = global_r
            else:
                r = np.random.rand()

            nm = eta_mut + 1.0

            # Calcular delta_q segun el valor de r
            if r < 0.5:
                bl = 2.0 * r + (1.0 - 2.0 * r) * ((1.0 - d) ** nm)
                delta_q = (bl ** (1.0 / nm)) - 1.0
            else:
                bl = 2.0 * (1.0 - r) + 2.0 * (r - 0.5) * ((1.0 - d) ** nm)
                delta_q = 1.0 - (bl ** (1.0 / nm))
            
            # Calcular la nueva posicion y asegurarse que este dentro de los limites
            y = x + delta_q * (xu - xl)
            mutant[i] = np.clip(y, xl, xu)
    
    return mutant
\end{lstlisting}

\section{Archivo \texttt{auxiliares\_functions.py}}
\begin{lstlisting}[
  caption={Implementación de \texttt{auxiliares\_functions.py}},
  style=pythonstyle,
  basicstyle=\ttfamily\footnotesize
]
import numpy as np
# ---------------------------
# Funciones auxiliares del GA
# ---------------------------
def initialize_population(pop_size, num_variables, lower_bound, upper_bound):
    """Inicializa la poblacion uniformemente en el espacio de busqueda."""
    return np.random.uniform(low=lower_bound, high=upper_bound, size=(pop_size, num_variables))
\end{lstlisting}
%\chapter{Registro de indicadores completo.}

\section{Langermann}

\subsection{Resumenes}
\begin{longtable}{lrrr}
\caption{Archivo: resumen\_run\_1}\label{tab:resumen_run_1} \\
\toprule
Indicador & x1 & x2 & Fitness \\
\midrule
\endfirsthead
\toprule
Indicador & x1 & x2 & Fitness \\
\midrule
\endhead
\midrule
\multicolumn{4}{r}{Continued on next page} \\
\midrule
\endfoot
\bottomrule
\endlastfoot
Mejor & 2.095881924190502 & 0.7727067906304386 & -4.914015915528873 \\
Media & 2.153832153948846 & 0.6853957759826199 & 1.2986960914753678 \\
Peor & 1.3636255150631562 & 0.2468636587424967 & 3.6320231404629033 \\
Desv. estándar & nan & nan & 2.442450990402752 \\
\end{longtable}
\begin{longtable}{lrrr}
\caption{Archivo: resumen\_run\_2}\label{tab:resumen_run_2} \\
\toprule
Indicador & x1 & x2 & Fitness \\
\midrule
\endfirsthead
\toprule
Indicador & x1 & x2 & Fitness \\
\midrule
\endhead
\midrule
\multicolumn{4}{r}{Continued on next page} \\
\midrule
\endfoot
\bottomrule
\endlastfoot
Mejor & -0.0002931428599729 & -0.000258734700225 & -0.9999944582414804 \\
Media & 0.1595412824455356 & 0.1595757109591133 & -0.4587742090713977 \\
Peor & -0.1835011970216257 & -0.1834667508654865 & -0.000189539565206 \\
Desv. estándar & nan & nan & 0.4531946247143622 \\
\end{longtable}
\begin{longtable}{lrrr}
\caption{Archivo: resumen\_run\_3}\label{tab:resumen_run_3} \\
\toprule
Indicador & x1 & x2 & Fitness \\
\midrule
\endfirsthead
\toprule
Indicador & x1 & x2 & Fitness \\
\midrule
\endhead
\midrule
\multicolumn{4}{r}{Continued on next page} \\
\midrule
\endfoot
\bottomrule
\endlastfoot
Mejor & 2.0167622431775043 & 0.9775688923221356 & -5.14581328834533 \\
Media & 0.8578923228636298 & 0.1119520311959549 & -2.034501886600382 \\
Peor & 1.158397366850342 & 0.1061677488445461 & -0.0597089267754305 \\
Desv. estándar & nan & nan & 0.5028469837532046 \\
\end{longtable}
\begin{longtable}{lrrr}
\caption{Archivo: resumen\_run\_4}\label{tab:resumen_run_4} \\
\toprule
Indicador & x1 & x2 & Fitness \\
\midrule
\endfirsthead
\toprule
Indicador & x1 & x2 & Fitness \\
\midrule
\endhead
\midrule
\multicolumn{4}{r}{Continued on next page} \\
\midrule
\endfoot
\bottomrule
\endlastfoot
Mejor & 0.0053825299661962 & -0.0039993874425109 & -0.9983708477155168 \\
Media & -0.611783101790296 & -0.6212321662131294 & -0.5334429338289646 \\
Peor & -0.9538970818830412 & -0.9634085976983588 & -0.0525302321903459 \\
Desv. estándar & nan & nan & 0.3296771206083393 \\
\end{longtable}
\begin{longtable}{lrrr}
\caption{Archivo: resumen\_run\_5}\label{tab:resumen_run_5} \\
\toprule
Indicador & x1 & x2 & Fitness \\
\midrule
\endfirsthead
\toprule
Indicador & x1 & x2 & Fitness \\
\midrule
\endhead
\midrule
\multicolumn{4}{r}{Continued on next page} \\
\midrule
\endfoot
\bottomrule
\endlastfoot
Mejor & -0.0374654803182066 & 0.5188910412092549 & -0.9362453043815636 \\
Media & -1.6294268221676989 & -1.0699856156682803 & -0.2407101556870587 \\
Peor & -1.1641102614113965 & -0.5983161207857252 & -4.204333777682678e-07 \\
Desv. estándar & nan & nan & 0.2149759772679994 \\
\end{longtable}
\begin{longtable}{lrrr}
\caption{Archivo: resumen\_run\_6}\label{tab:resumen_run_6} \\
\toprule
Indicador & x1 & x2 & Fitness \\
\midrule
\endfirsthead
\toprule
Indicador & x1 & x2 & Fitness \\
\midrule
\endhead
\midrule
\multicolumn{4}{r}{Continued on next page} \\
\midrule
\endfoot
\bottomrule
\endlastfoot
Mejor & -0.0120281638507773 & 0.0119344030286503 & -0.9896286873304896 \\
Media & 1.1078639459141058 & 1.1319204270442176 & -0.5447942472082268 \\
Peor & 1.1458796638728408 & 1.169938981286806 & -0.5074559004809934 \\
Desv. estándar & nan & nan & 0.0622931458098346 \\
\end{longtable}
\begin{longtable}{lrrr}
\caption{Archivo: resumen\_run\_7}\label{tab:resumen_run_7} \\
\toprule
Indicador & x1 & x2 & Fitness \\
\midrule
\endfirsthead
\toprule
Indicador & x1 & x2 & Fitness \\
\midrule
\endhead
\midrule
\multicolumn{4}{r}{Continued on next page} \\
\midrule
\endfoot
\bottomrule
\endlastfoot
Mejor & 2.0019436362996443 & 1.0037506288328364 & -5.162016654905215 \\
Media & 0.1192414512066273 & 0.0353625964394437 & 0.6120035661402777 \\
Peor & 0.0252458179930767 & 0.003505827249705 & 1.0031083996798786 \\
Desv. estándar & nan & nan & 0.9666580436507198 \\
\end{longtable}
\begin{longtable}{lrrr}
\caption{Archivo: resumen\_run\_8}\label{tab:resumen_run_8} \\
\toprule
Indicador & x1 & x2 & Fitness \\
\midrule
\endfirsthead
\toprule
Indicador & x1 & x2 & Fitness \\
\midrule
\endhead
\midrule
\multicolumn{4}{r}{Continued on next page} \\
\midrule
\endfoot
\bottomrule
\endlastfoot
Mejor & 0.0153481319147759 & -0.0152376917924069 & -0.9831402532063 \\
Media & -1.906644597007218 & -1.9384580477844384 & -0.2019478608654046 \\
Peor & -1.9878786813725344 & -2.019433856241915 & -0.0248051240544462 \\
Desv. estándar & nan & nan & 0.1303563220305934 \\
\end{longtable}
\begin{longtable}{lrrr}
\caption{Archivo: resumen\_run\_9}\label{tab:resumen_run_9} \\
\toprule
Indicador & x1 & x2 & Fitness \\
\midrule
\endfirsthead
\toprule
Indicador & x1 & x2 & Fitness \\
\midrule
\endhead
\midrule
\multicolumn{4}{r}{Continued on next page} \\
\midrule
\endfoot
\bottomrule
\endlastfoot
Mejor & 2.002592498731848 & 1.006380979616417 & -5.162121824564412 \\
Media & 2.1147866299358484 & 1.1024062009290632 & -4.895316803411505 \\
Peor & 2.147627653713365 & 1.132626689731959 & -4.804433929172063 \\
Desv. estándar & nan & nan & 0.0509847493971701 \\
\end{longtable}
\begin{longtable}{lrrr}
\caption{Archivo: resumen\_run\_10}\label{tab:resumen_run_10} \\
\toprule
Indicador & x1 & x2 & Fitness \\
\midrule
\endfirsthead
\toprule
Indicador & x1 & x2 & Fitness \\
\midrule
\endhead
\midrule
\multicolumn{4}{r}{Continued on next page} \\
\midrule
\endfoot
\bottomrule
\endlastfoot
Mejor & 2.0058434821635696 & 1.0004095459845868 & -5.16143210203958 \\
Media & 0.8724703163431755 & 0.1101092286402504 & -2.377918278148532 \\
Peor & 1.3775297218210725 & 0.4582047123613642 & 2.187761548577683 \\
Desv. estándar & nan & nan & 0.4930068965724287 \\
\end{longtable}

\section{Drop Wave}

\subsection{Resumenes}
\begin{longtable}{lrrr}
\caption{Archivo: resumen\_run\_1}\label{tab:resumen_run_1} \\
\toprule
Indicador & x1 & x2 & Fitness \\
\midrule
\endfirsthead
\toprule
Indicador & x1 & x2 & Fitness \\
\midrule
\endhead
\midrule
\multicolumn{4}{r}{Continued on next page} \\
\midrule
\endfoot
\bottomrule
\endlastfoot
Mejor & 2.095881924190502 & 0.7727067906304386 & -4.914015915528873 \\
Media & 2.153832153948846 & 0.6853957759826199 & 1.2986960914753678 \\
Peor & 1.3636255150631562 & 0.2468636587424967 & 3.6320231404629033 \\
Desv. estándar & nan & nan & 2.442450990402752 \\
\end{longtable}
\begin{longtable}{lrrr}
\caption{Archivo: resumen\_run\_2}\label{tab:resumen_run_2} \\
\toprule
Indicador & x1 & x2 & Fitness \\
\midrule
\endfirsthead
\toprule
Indicador & x1 & x2 & Fitness \\
\midrule
\endhead
\midrule
\multicolumn{4}{r}{Continued on next page} \\
\midrule
\endfoot
\bottomrule
\endlastfoot
Mejor & -0.0002931428599729 & -0.000258734700225 & -0.9999944582414804 \\
Media & 0.1595412824455356 & 0.1595757109591133 & -0.4587742090713977 \\
Peor & -0.1835011970216257 & -0.1834667508654865 & -0.000189539565206 \\
Desv. estándar & nan & nan & 0.4531946247143622 \\
\end{longtable}
\begin{longtable}{lrrr}
\caption{Archivo: resumen\_run\_3}\label{tab:resumen_run_3} \\
\toprule
Indicador & x1 & x2 & Fitness \\
\midrule
\endfirsthead
\toprule
Indicador & x1 & x2 & Fitness \\
\midrule
\endhead
\midrule
\multicolumn{4}{r}{Continued on next page} \\
\midrule
\endfoot
\bottomrule
\endlastfoot
Mejor & 2.0167622431775043 & 0.9775688923221356 & -5.14581328834533 \\
Media & 0.8578923228636298 & 0.1119520311959549 & -2.034501886600382 \\
Peor & 1.158397366850342 & 0.1061677488445461 & -0.0597089267754305 \\
Desv. estándar & nan & nan & 0.5028469837532046 \\
\end{longtable}
\begin{longtable}{lrrr}
\caption{Archivo: resumen\_run\_4}\label{tab:resumen_run_4} \\
\toprule
Indicador & x1 & x2 & Fitness \\
\midrule
\endfirsthead
\toprule
Indicador & x1 & x2 & Fitness \\
\midrule
\endhead
\midrule
\multicolumn{4}{r}{Continued on next page} \\
\midrule
\endfoot
\bottomrule
\endlastfoot
Mejor & 0.0053825299661962 & -0.0039993874425109 & -0.9983708477155168 \\
Media & -0.611783101790296 & -0.6212321662131294 & -0.5334429338289646 \\
Peor & -0.9538970818830412 & -0.9634085976983588 & -0.0525302321903459 \\
Desv. estándar & nan & nan & 0.3296771206083393 \\
\end{longtable}
\begin{longtable}{lrrr}
\caption{Archivo: resumen\_run\_5}\label{tab:resumen_run_5} \\
\toprule
Indicador & x1 & x2 & Fitness \\
\midrule
\endfirsthead
\toprule
Indicador & x1 & x2 & Fitness \\
\midrule
\endhead
\midrule
\multicolumn{4}{r}{Continued on next page} \\
\midrule
\endfoot
\bottomrule
\endlastfoot
Mejor & -0.0374654803182066 & 0.5188910412092549 & -0.9362453043815636 \\
Media & -1.6294268221676989 & -1.0699856156682803 & -0.2407101556870587 \\
Peor & -1.1641102614113965 & -0.5983161207857252 & -4.204333777682678e-07 \\
Desv. estándar & nan & nan & 0.2149759772679994 \\
\end{longtable}
\begin{longtable}{lrrr}
\caption{Archivo: resumen\_run\_6}\label{tab:resumen_run_6} \\
\toprule
Indicador & x1 & x2 & Fitness \\
\midrule
\endfirsthead
\toprule
Indicador & x1 & x2 & Fitness \\
\midrule
\endhead
\midrule
\multicolumn{4}{r}{Continued on next page} \\
\midrule
\endfoot
\bottomrule
\endlastfoot
Mejor & -0.0120281638507773 & 0.0119344030286503 & -0.9896286873304896 \\
Media & 1.1078639459141058 & 1.1319204270442176 & -0.5447942472082268 \\
Peor & 1.1458796638728408 & 1.169938981286806 & -0.5074559004809934 \\
Desv. estándar & nan & nan & 0.0622931458098346 \\
\end{longtable}
\begin{longtable}{lrrr}
\caption{Archivo: resumen\_run\_7}\label{tab:resumen_run_7} \\
\toprule
Indicador & x1 & x2 & Fitness \\
\midrule
\endfirsthead
\toprule
Indicador & x1 & x2 & Fitness \\
\midrule
\endhead
\midrule
\multicolumn{4}{r}{Continued on next page} \\
\midrule
\endfoot
\bottomrule
\endlastfoot
Mejor & 2.0019436362996443 & 1.0037506288328364 & -5.162016654905215 \\
Media & 0.1192414512066273 & 0.0353625964394437 & 0.6120035661402777 \\
Peor & 0.0252458179930767 & 0.003505827249705 & 1.0031083996798786 \\
Desv. estándar & nan & nan & 0.9666580436507198 \\
\end{longtable}
\begin{longtable}{lrrr}
\caption{Archivo: resumen\_run\_8}\label{tab:resumen_run_8} \\
\toprule
Indicador & x1 & x2 & Fitness \\
\midrule
\endfirsthead
\toprule
Indicador & x1 & x2 & Fitness \\
\midrule
\endhead
\midrule
\multicolumn{4}{r}{Continued on next page} \\
\midrule
\endfoot
\bottomrule
\endlastfoot
Mejor & 0.0153481319147759 & -0.0152376917924069 & -0.9831402532063 \\
Media & -1.906644597007218 & -1.9384580477844384 & -0.2019478608654046 \\
Peor & -1.9878786813725344 & -2.019433856241915 & -0.0248051240544462 \\
Desv. estándar & nan & nan & 0.1303563220305934 \\
\end{longtable}
\begin{longtable}{lrrr}
\caption{Archivo: resumen\_run\_9}\label{tab:resumen_run_9} \\
\toprule
Indicador & x1 & x2 & Fitness \\
\midrule
\endfirsthead
\toprule
Indicador & x1 & x2 & Fitness \\
\midrule
\endhead
\midrule
\multicolumn{4}{r}{Continued on next page} \\
\midrule
\endfoot
\bottomrule
\endlastfoot
Mejor & 2.002592498731848 & 1.006380979616417 & -5.162121824564412 \\
Media & 2.1147866299358484 & 1.1024062009290632 & -4.895316803411505 \\
Peor & 2.147627653713365 & 1.132626689731959 & -4.804433929172063 \\
Desv. estándar & nan & nan & 0.0509847493971701 \\
\end{longtable}
\begin{longtable}{lrrr}
\caption{Archivo: resumen\_run\_10}\label{tab:resumen_run_10} \\
\toprule
Indicador & x1 & x2 & Fitness \\
\midrule
\endfirsthead
\toprule
Indicador & x1 & x2 & Fitness \\
\midrule
\endhead
\midrule
\multicolumn{4}{r}{Continued on next page} \\
\midrule
\endfoot
\bottomrule
\endlastfoot
Mejor & 2.0058434821635696 & 1.0004095459845868 & -5.16143210203958 \\
Media & 0.8724703163431755 & 0.1101092286402504 & -2.377918278148532 \\
Peor & 1.3775297218210725 & 0.4582047123613642 & 2.187761548577683 \\
Desv. estándar & nan & nan & 0.4930068965724287 \\
\end{longtable}

\end{document}


\end{document}
